% Build from the repository root with: bash scripts/build_cpsv17_corrections.sh
\documentclass[aps,pra,reprint,nofootinbib]{revtex4-2}
\usepackage[T1]{fontenc}
\usepackage{amsmath,amssymb,amsthm,bm,dsfont,graphicx,booktabs,microtype}
\usepackage{xr-hyper}
\usepackage[hidelinks]{hyperref}
\usepackage{cleveref}
\crefname{equation}{Eq.}{Eqs.}
\IfFileExists{../../build/paper-aux/1606.00608/MPDO-22-12-17-2-note.aux}{%
  \externaldocument[cpsv-]{../../build/paper-aux/1606.00608/MPDO-22-12-17-2-note}[https://arxiv.org/pdf/1606.00608]%
}{%
  \IfFileExists{build/paper-aux/1606.00608/MPDO-22-12-17-2-note.aux}{%
    \externaldocument[cpsv-]{build/paper-aux/1606.00608/MPDO-22-12-17-2-note}[https://arxiv.org/pdf/1606.00608]%
  }{\externaldocument[cpsv-]{MPDO-22-12-17-2-note}[https://arxiv.org/pdf/1606.00608]}%
}
\newcommand{\tr}{\operatorname{tr}}
\newcommand{\Id}{\mathds{1}}
\newcommand{\ket}[1]{\lvert#1\rangle}
\newcommand{\bra}[1]{\langle#1\rvert}
\newcommand{\ketbra}[2]{\lvert#1\rangle\langle#2\rvert}
\newcommand{\mc}[1]{\mathcal{#1}}
\newcommand{\diag}{\operatorname{diag}}
\newtheorem{proposition}{Proposition}
\newtheorem{lemma}{Lemma}
\setcounter{secnumdepth}{3}
\setlength{\emergencystretch}{1em}

\begin{document}
\title{Corrections to matrix product density operator fixed-point statements}
\begin{abstract}
	A simple matrix product density operator tensor in block-injective
	canonical form can have literal zero correlation length and exact
	finite-ring saturation of the area law without admitting the blocking
	channels asserted in Proposition~\ref{cpsv-prop2to5} of
	Cirac--P\'erez-Garc\'ia--Schuch--Verstraete.
	We give an explicit example satisfying the global canonical-weight
	normalization. A Hermitian boundary closure increases in trace norm
	from two sites to either three or four sites, excluding the required
	refinement channels and the implied two-site-blocked fixed point.
	We also give compact counterexamples involving canonical copy weights,
	peripheral phases, and finite-periodic entropy. These are separated from
	local corrections and from conclusions that remain valid, including
	the pure renormalization-flow limit characterization and active-sector
	matching in the fundamental theorem.
\end{abstract}
\maketitle

\section{Scope and conventions}\label{sec:scope}

Proposition~\ref{cpsv-prop2to5} of
Ref.~\cite{Cirac2017MPDO_AnnPhys} is false under all its printed standing
hypotheses. The example below also disproves the implication from
condition~(ii) to condition~(v) in the displayed chain of
Theorem~\ref{cpsv-thm:main-simple}.
Source references use the numbering of arXiv:1606.00608v4;
unqualified equation references are local.

For an MPDO tensor $K$, with physical dimension $d$ and bond dimension
$D$, define its boundary closure by
\begin{align}
	K_N(X)
	                & =\sum_{\boldsymbol i,\boldsymbol j}
	\tr\!\left(XK^{i_1j_1}\cdots K^{i_Nj_N}\right)
	\ket{\boldsymbol i}\bra{\boldsymbol j},
	\label{eq:closure}                                    \\
	\sigma^{(N)}(K) & =K_N(\Id_D).
	\nonumber
\end{align}
Here $X$ is any complex $D\times D$ matrix. Two different transfer
objects must be distinguished:
\begin{align*}
	B_K        & =\sum_iK^{ii},                      &
	\mc E_K(H) & =\sum_{i,j}K^{ij}H(K^{ij})^\dagger.
\end{align*}
Mixed-state zero correlation length (ZCL) is the literal identity
$B_K^2=B_K$. A renormalization fixed point (RFP) requires
trace-preserving completely positive maps $\mc T,\mc S$ with
$\mc T[K_1(X)]=K_2(X)$ and $\mc S[K_2(X)]=K_1(X)$ for every $X$.

We use the source definitions of canonical form (CF), normal tensors,
a basis of normal tensors (BNT), and block-injective canonical form (biCF).
Normal representatives are irreducible and have completely positive
transfer maps of spectral radius one, with no other peripheral
eigenvalue. In biCF their one-site tuples span the direct sum of the
representatives' full matrix algebras.
Canonical weights obey $|\mu_{j,q}|\leq1$, with at least one
unit-modulus weight globally, not necessarily in every sector.
Simplicity means that no BNT representative has nilpotent
physical-trace matrix.

Entropies concern normalized periodic operators and use base-two
logarithms. For a connected block of $L$ sites, write
\[
	I_L^{(N)}=S_L^{(N)}+S_{N-L}^{(N)}-S_N^{(N)}.
\]
Saturation of the area law (SAL) means equality of these mutual
informations up to half the ring. Normalization for entropy calculations
does not rescale the tensor in a channel identity.

\section{Saturation and zero correlation length do not imply the stated blocking channel}\label{sec:blocking}

\subsection{The four-letter tensor}

Let $i=0,1,2,3$, and set
\begin{align*}
	t   & =(1,2,-7,4)^{\mathsf T},                  &
	s   & =(-3,-1,1,3)^{\mathsf T},                   \\
	l_i & =\begin{pmatrix}1\\t_i\end{pmatrix},      &
	r_i & =\begin{pmatrix}1/4&s_i/100\end{pmatrix}.
\end{align*}
Define a diagonal physical tensor $F$ by $F^{ii}=F_i=l_ir_i$,
with all off-diagonal physical slices zero. Its physical dimension
is four and its bond dimension is two. Let $L$ have columns $l_i$
and $R$ have rows $r_i$. The identities $\sum_i t_i=\sum_i s_i=\sum_i s_it_i=0$ give
\begin{align}
	LR & =P=\diag(1,0), \label{eq:LR} \\
	T:=RL
	   & =\frac1{100}
	\begin{pmatrix}
		22 & 19 & 46 & 13 \\
		24 & 23 & 32 & 21 \\
		26 & 27 & 18 & 29 \\
		28 & 31 & 4  & 37
	\end{pmatrix}.
	\label{eq:T}
\end{align}
In particular, $T_{ij}=1/4+s_it_j/100$ is strictly positive,
and every row and column sums to one.

Let $\boldsymbol1=(1,1,1,1)^{\mathsf T}$. Put
$Q=\boldsymbol1\boldsymbol1^{\mathsf T}/4$ and
$N_0=st^{\mathsf T}/100$. The three scalar sums above imply
\[
	Q^2=Q,\qquad QN_0=N_0Q=N_0^2=0.
\]
Since $T=Q+N_0$, we have the exact identities
\begin{align}
	T^2 & =Q,\qquad T^m=Q\quad(m\geq2).
	\label{eq:Tpowers}
\end{align}
Rank-one multiplication also gives
\begin{align}
	F_{i_1}\cdots F_{i_N}
	 & =l_{i_1}\left(\prod_{a=1}^{N-1}T_{i_ai_{a+1}}\right)r_{i_N}.
	\label{eq:rankone}
\end{align}
Its trace is $\prod_{a=1}^{N}T_{i_ai_{a+1}}$, with
$i_{N+1}=i_1$. Thus every periodic operator is positive and
$\tr\sigma^{(N)}(F)=\tr T^N=1$, since $\tr T=\tr Q=1$.
Moreover $B_F=LR=P$, so $B_F^2=B_F$.
Neither positivity nor literal ZCL requires a rescaling.

\subsection{Exact finite-ring SAL}

The normalized periodic distribution is
$p_N(i_1,\ldots,i_N)=\prod_aT_{i_ai_{a+1}}$.
For every proper connected block, summing over the complementary
path gives
\begin{align}
	p_{N,L}(i_1,\ldots,i_L)
	 & =\left(\prod_{a=1}^{L-1}T_{i_ai_{a+1}}\right)
	    (T^{N-L+1})_{i_Li_1}\nonumber \\
	 & =\frac14\prod_{a=1}^{L-1}T_{i_ai_{a+1}},
	\qquad 1\leq L<N.
	\label{eq:marginal}
\end{align}
The complementary exponent is at least two, so
\cref{eq:Tpowers} applies even when only one site is traced out.

Define the mean row entropy
\[
	h=-\frac14\sum_{i,j}T_{ij}\log_2T_{ij}.
\]
The marginal in \cref{eq:marginal} has a uniform initial letter
and internal edge probabilities $T_{ij}/4$. Therefore
\begin{align*}
	S_L^{(N)}
	 & =-\mathbb E\!\left[-2+\sum_{a=1}^{L-1}\log_2T_{i_ai_{a+1}}\right] \\
	 & =2+(L-1)h.
\end{align*}
For $N\geq3$, each edge of the full ring also has distribution
$T_{ij}(T^{N-1})_{ji}=T_{ij}/4$. Hence
\begin{align*}
	S_N^{(N)} & =Nh,                    \\
	I_L^{(N)}
	          & =2+(L-1)h+2+(N-L-1)h-Nh \\
	          & =4-2h.
\end{align*}
This proves exact finite-ring SAL. At $N=2$ there is only one
nonempty cut size and thus no additional equality to check.

\subsection{Normal form and the global canonical weight}

To meet the standing canonical normalization, the four-letter
tensor must be placed inside a slightly larger tensor rather than
rescaled arbitrarily. First, its slices span $M_2(\mathbb C)$.
Their coordinate rows, in the order
$(11,12,21,22)$, are
$(1/4,s_i/100,t_i/4,s_it_i/100)$.
The determinant of this four-by-four coordinate matrix is
$-3/2000\ne0$, proving the span assertion.

This span makes $\mc E_F$ strictly positive. Indeed, for
$H\geq0$, $H\ne0$, and $v\ne0$,
\[
	v^\dagger\mc E_F(H)v
	=\sum_i\|H^{1/2}F_i^\dagger v\|^2>0.
\]
Equality would imply $H^{1/2}A^\dagger v=0$ for every
$A\in M_2(\mathbb C)$, although these vectors $A^\dagger v$
span $\mathbb C^2$. Finite-dimensional Perron--Frobenius theory
then supplies a Perron value $r>0$, positive-definite left and
right eigenmatrices, and no other peripheral eigenvalue.

We also need $r<1$. Write $\mc E_F^*$ for the trace adjoint.
For $H=\diag(20,1)$, direct substitution gives
\[
	H-\mc E_F^*(H)
	=\begin{pmatrix}85/8&-9/40\\-9/40&4697/5000\end{pmatrix}>0,
\]
Its leading principal minor is positive and its determinant is
$19861/2000>0$.
Pairing with a right Perron eigenmatrix $\rho>0$ gives
\begin{align*}
	0 & <\tr[(H-\mc E_F^*(H))\rho]      \\
	  & =\tr(H\rho)-\tr(H\mc E_F(\rho)) \\
	  & =(1-r)\tr(H\rho),
\end{align*}
so $0<r<1$.

Let $\mu=\sqrt r$. If $H_*>0$ is a left Perron eigenmatrix,
set $G=H_*^{1/2}$ and
$\mc K_0^{ii}=\mu^{-1}GF_iG^{-1}$. Then
\[
	\sum_i(\mc K_0^{ii})^\dagger\mc K_0^{ii}
	=r^{-1}G^{-1}\mc E_F^*(H_*)G^{-1}=\Id_2.
\]
Thus $\mc K_0$ is a normal representative with $0<\mu<1$.

Embed its physical support into the first four coordinates of
$\mathbb C^5$. On the fifth coordinate take the bond-one normal
tensor $\mc K_1^{44}=1$, with all other slices zero. Define
\begin{align}
	K=(\mu\mc K_0)\oplus\mc K_1.
	\label{eq:ambient}
\end{align}
This tensor has physical dimension five and bond dimension three.
The representatives have orthogonal physical supports and simultaneous
one-site span $M_2(\mathbb C)\oplus\mathbb C$, proving biCF.
Their weights $\mu,1$ satisfy the printed global normalization.
Their physical-trace matrices are
$\mu^{-1}GPG^{-1}$ and $1$, neither nilpotent; hence $K$ is simple.

The periodic operator is the orthogonal sum of the embedded
$\sigma^{(N)}(F)$ and $\ketbra{4^N}{4^N}$, both of trace one.
It is positive with trace two at every positive length.
Every nonempty marginal, including the full state, satisfies
\[
	S_L^{(N)}(K)=1+\tfrac12S_L^{(N)}(F).
\]
It follows that
$I_L^{(N)}(K)=1+\frac12I_L^{(N)}(F)$, proving SAL.
Finally,
\[
	B_K=(GPG^{-1})\oplus1,\qquad B_K^2=B_K.
\]
All standing hypotheses hold simultaneously.

\subsection{A boundary excluding every refinement channel}

Take
\[
	X=\begin{pmatrix}-1&2/5\\-5&0\end{pmatrix}.
\]
The endpoint coefficients derived from \cref{eq:rankone} are
\begin{align*}
	D_{ij}=r_jXl_i
	  & =-\frac14+\frac{t_i}{10}-\frac{s_j}{20}, \\
	D & =\begin{pmatrix}
		     0    & -1/10 & -1/5  & -3/10  \\
		     1/10 & 0     & -1/10 & -1/5   \\
		     -4/5 & -9/10 & -1    & -11/10 \\
		     3/10 & 1/5   & 1/10  & 0
	     \end{pmatrix}.
\end{align*}
Every $F_N(X)$ is real diagonal, with coefficient
$D_{i_1i_N}\prod_{a=1}^{N-1}T_{i_ai_{a+1}}$.
It is therefore Hermitian, even though $X$ need not be.
Positivity of the path weights gives
\begin{align}
	\|F_N(X)\|_1
	=\sum_{i,j}|D_{ij}|(T^{N-1})_{ij}.
	\label{eq:norm}
\end{align}
Using the displayed matrices and \cref{eq:Tpowers} gives
\begin{align}
	\|F_2(X)\|_1              & =\frac{337}{250},\nonumber \\
	\|F_3(X)\|_1=\|F_4(X)\|_1 & =\frac{27}{20}.
	\label{eq:norm-gap}
\end{align}
The increase is $27/20-337/250=1/500>0$.

Let $V:\mathbb C^4\longrightarrow\mathbb C^5$ be the coordinate
inclusion. Use the same ambient boundary
$Y=(GXG^{-1})\oplus0$ at every length. Cyclicity of the trace
cancels the virtual gauges, while the zero terminal boundary removes
the fifth-letter sector:
\begin{align}
	K_N(Y)=V^{\otimes N}F_N(X)(V^{\otimes N})^\dagger.
	\label{eq:boundary}
\end{align}
This isometric embedding preserves all singular values except for
adding zeros, so \cref{eq:norm-gap} holds unchanged.

A positive trace-preserving map $\Phi$ contracts the trace norm of
Hermitian matrices. In fact, for $Z=Z_+-Z_-$,
\begin{align*}
	\|\Phi(Z)\|_1
	 & \leq\|\Phi(Z_+)\|_1+\|\Phi(Z_-)\|_1 \\
	 & =\tr\Phi(Z_+)+\tr\Phi(Z_-)=\|Z\|_1.
\end{align*}
Thus no channel can send $K_2(Y)$ to $K_3(Y)$ or to $K_4(Y)$.

Proposition~\ref{cpsv-prop2to5}, which requires the two-to-three-site
identity for every boundary, is false. The two-site-blocked tensor
in Theorem~\ref{cpsv-thm:main-simple} is not an RFP either, because
its refinement would send $K_2(Y)$ to $K_4(Y)$.
This excludes all channels, not merely the channel obtained from
one proposed factorization. It makes no assertion about every larger
blocking and does not refute Theorem~\ref{cpsv-thm:IV.13}.

\subsection{Two points in the structural proof}

The example also identifies two distinct problems in the proposed
route to the channel. First, primitivity and traces of powers do
not eliminate zero-eigenvalue Jordan blocks. The strictly positive
matrix $T=Q+N_0$ above has one eigenvalue one and all other
eigenvalues zero, yet is not rank one:
\[
	T_{00}T_{11}-T_{01}T_{10}
	=\frac{22\cdot23-19\cdot24}{10000}=\frac1{200}.
\]
The nonzero nilpotent term $N_0^2=0$ is invisible in
$\tr T^N$. This invalidates the rank-one inference in source
Lemma~\ref{cpsv-SALZCL}.

Second, common-weight absorption after source
Lemma~\ref{cpsv-lemmus} does not generally preserve normality:
\[
	\mc E_{\mu_j\mc K_j}=|\mu_j|^2\mc E_{\mc K_j}.
\]
A normal representative remains normal after this absorption only
if $|\mu_j|=1$. In \cref{eq:ambient}, the first absorbed block
has spectral radius $r<1$.
These observations explain defects in the structural argument.
The channel refutation above does not require a further universal
claim about all possible physical-sector factorizations.

\section{Raw weights and peripheral phases}\label{sec:weights-phases}

\subsection{Physical ZCL does not determine raw weights}

For pure tensors, the source defines ZCL as local orthogonality of the
BNT representatives together with correlations independent of distance
(CID). This is a physical condition, distinct from the mixed-state
transfer identity in Sec.~\ref{sec:scope}.
For a pure tensor $A$, write
$\mc E_A(H)=\sum_iA^iH(A^i)^\dagger$ and
$\mathbb E_A=\sum_iA^i\otimes\overline{A^i}$ for its transfer map
and matrix. In canonical coordinates, the copy block $(j,q;j',q')$ is
\[
	\mu_{j,q}\overline{\mu_{j',q'}}
	\sum_iA_j^i\otimes\overline{A_{j'}^i}.
\]
The copy factors matter for literal idempotence.

Take one physical letter and $A^0=\diag(1,1/2)$.
This is in CF with one scalar normal BNT representative, two
copies, and allowed weights $1,1/2$. Its periodic vector is
\[
	\ket{V^{(N)}(A)}=(1+2^{-N})\ket{0^N}.
\]
Local orthogonality is vacuous. Every physical observable is scalar,
so correlations are independent of separation, even for unnormalized
expectations at fixed $N$. Thus the printed physical ZCL condition
holds. But for the matrix unit $e_{12}$,
\[
	\mc E_A(e_{12})=\tfrac12e_{12},\qquad
	\mc E_A^2(e_{12})=\tfrac14e_{12}.
\]
This refutes the ZCL-to-idempotence direction of source
Theorem~\ref{cpsv-TheoremZCLPure}, and hence the unrestricted
equivalence in Theorem~\ref{cpsv-thm:main-MPS}.
Directly, the one-letter RFP equation would require
$(A^0)^2=uA^0$, whose diagonal entries force $u=1$ and
$1/4=1/2$.

Regard the same matrix as an MPDO tensor
$K^{00}=\diag(1,1/2)$. It gives an independent counterexample to
Theorem~\ref{cpsv-thm:main-simple}(iv)$\Rightarrow$(v).
It is simple and in biCF, has the required global weight normalization,
and generates the positive scalar $1+2^{-N}$. Its sole normal
representative is $1$, whose scalar physical-sector factorization
has $\eta=1$ and $a=b=1$. Thus condition~(iv) holds in full.
Nevertheless, trace preservation for the two-site-blocked tensor
would require
\[
	(K^{00})^4=(K^{00})^2,
	\qquad
	\diag(1,1/16)=\diag(1,1/4),
\]
which is false. Condition~(iv) constrains normal representatives,
not their raw copy weights. This independent example does not
satisfy mixed-state literal ZCL and is not the counterexample to
Proposition~\ref{cpsv-prop2to5}.

\subsection{Peripheral phases and similarity under squaring}

Let
\[
	\omega=e^{2\pi i/3},\qquad
	A^0=\diag(1,\omega),\qquad
	S=\begin{pmatrix}0&1\\1&0\end{pmatrix}.
\]
Both scalar normal copies have unit-modulus weights. Nevertheless,
\[
	\mc E_A^{2^n}(e_{21})=\omega^{2^n}e_{21}.
\]
Since $2^{2m}\equiv1$ and $2^{2m+1}\equiv2\pmod3$, this
sequence alternates between $\omega e_{21}$ and
$\omega^2e_{21}$. The universal convergence assertion following
source Theorem~\ref{cpsv-thm:renormalization-flow} and in
Appendix~B is therefore false.

The same example disproves the pure-state equivalence asserted after
source Eq.~\eqref{cpsv-RFP-gauge}. Indeed,
\begin{align*}
	\omega S(A^0)^2S^{-1}
	 & =\omega\diag(\omega^2,1)=A^0.
\end{align*}
For the pure MPDO tensor $M=A^0\otimes\overline{A^0}$ the phase
cancels, so
\[
	M=(S\otimes\overline S)M^2(S\otimes\overline S)^{-1}.
\]
The one-dimensional physical algebras have identity channels, giving
the source's blocking-up-to-virtual-gauge relations in both directions.
The doubled tensor is itself in CF with scalar normal copies and
weights $1,\overline\omega,\omega,1$, and generates the positive
scalar $|1+\omega^N|^2$. But $M^2\ne M$, so it is not a literal
channel RFP. Similarity under squaring permits nontrivial
root-of-unity cycles; it does not force spectrum in $\{0,1\}$.

Source Theorem~\ref{cpsv-thm:renormalization-flow} itself survives
as a limit characterization. If $\mc E^{2^n}\to F$, continuity
of composition gives
\[
	F^2=\lim_n(\mc E^{2^n})^2
	=\lim_n\mc E^{2^{n+1}}=F.
\]
Conversely an idempotent completely positive map has a constant
squaring orbit. The correspondence between such maps and physical
isometry classes of Kraus tensors gives the asserted fixed-point
characterization. Existence of a limit for every CF initial tensor
is a separate, false assertion.

\section{Finite-chain entropy and thermodynamic limits}\label{sec:entropy}

\subsection{The inner thermodynamic limit}

Consider
\begin{align*}
	M^{00} & =\diag(1,0,0),  &
	M^{11} & =\diag(0,1,-1),                \\
	M^{01} & =0,             & M^{10} & =0.
\end{align*}
Mixed physical words vanish; constant words give
\begin{align*}
	\rho^{(N)}
	              & =\ketbra{0^N}{0^N}
	+(1+(-1)^N)\ketbra{1^N}{1^N},      \\
	\tr\rho^{(N)} & =2+(-1)^N.
\end{align*}
Every positive-length operator is positive. The tensor is in CF
with scalar normal copies and globally normalized weights $1,1,-1$.

At odd $N$ the normalized state and all nonempty marginals are
pure. At even $N$ they all have eigenvalues $1/3,2/3$.
For every fixed $L\geq1$ and $N>L$,
\[
	I_L^{(N)}
	=\begin{cases}
		0,        & N\text{ odd},  \\
		h_2(1/3), & N\text{ even},
	\end{cases}
\]
where $h_2(p)=-p\log_2p-(1-p)\log_2(1-p)$.
Since $h_2(1/3)>0$, the unrestricted inner limit in source
Proposition~\ref{cpsv-PropILILp1} does not exist.

Finite-chain monotonicity remains valid for
$1\leq L<\lfloor N/2\rfloor$. Strong subadditivity on
consecutive regions of lengths $1,L,N-2L-1$, together with
translation invariance, gives
\begin{align*}
	S_L^{(N)}+S_{N-L}^{(N)}
	          & \leq S_{L+1}^{(N)}+S_{N-L-1}^{(N)}, \\
	I_L^{(N)} & \leq I_{L+1}^{(N)}.
\end{align*}
The finite-chain bound $I_L^{(N)}\leq4\log_2D$ quoted in the
source proof is likewise separate from convergence.
Monotonicity in $L$ and a uniform bound do not establish a
limit in $N$. If the fixed-$L$ limits exist, these finite-chain
properties do imply that their subsequent limit as $L\to\infty$
exists and is bounded.

\subsection{The literal periodic state in Example~\ref{cpsv-ExSALnoZCL}}

The neighboring operators in source
Example~\ref{cpsv-ExSALnoZCL} place site $n$ in
$\ket{k_n,k_n}$. Its nonzero spectrum is therefore the classical
periodic distribution
\begin{align*}
	p_N(k_1,\ldots,k_N)
	  & =\frac1{Z_N}\prod_{n=1}^NW_{k_nk_{n+1}},   \\
	W & =\begin{pmatrix}1&1/2\\1/2&1\end{pmatrix},
	\qquad k_{N+1}=k_1.
\end{align*}
The spectral projectors
$P_\pm=\frac12(\begin{smallmatrix}1&\pm1\\\pm1&1\end{smallmatrix})$
give
\begin{align*}
	W^m & =(3/2)^mP_++(1/2)^mP_-,               \\
	Z_N & =\tr W^N=(3^N+1)/2^N,                 \\
	p_{N,L}(k_1,\ldots,k_L)
	    & =\frac{\prod_{n=1}^{L-1}W_{k_nk_{n+1}}
			       (W^{N-L+1})_{k_Lk_1}}{Z_N}.
\end{align*}
The complementary path cannot be omitted on a finite ring.

At $N=4$, $Z_4=41/8$, and
\[
	W^2=\begin{pmatrix}5/4&1\\1&5/4\end{pmatrix},
	\qquad
	W^3=\begin{pmatrix}7/4&13/8\\13/8&7/4\end{pmatrix}.
\]
The two equal-bit pair probabilities are
$(7/4)/(41/8)=14/41$; the two unequal-bit pair probabilities
are $(1/2)(13/8)/(41/8)=13/82$.
For three bits, zero, one, and two internal transitions give
probabilities $10/41,4/41,5/82$, with multiplicities $2,4,2$.
The one-bit marginal is uniform. Evaluating the entropies of these
probabilities in $I_2-I_1=2S_2-S_1-S_3$ gives
\[
	I_2-I_1
	=\frac1{82}\log_2
	\frac{41^{82}5^{50}}{2^{48}13^{52}7^{112}}.
\]

Strict positivity admits a short proof without decimal evaluation
or a large integer comparison. Translation invariance identifies
this defect with $I(k_1:k_3\mid k_2)$.
Conditioned on $k_2=0$, the joint table of $k_1,k_3$ is
\[
	q=\frac1{41}\begin{pmatrix}20&8\\8&5\end{pmatrix}.
\]
It is not a product distribution, since
$q_{00}q_{11}-q_{01}q_{10}=(100-64)/41^2>0$.
For positive probability distributions $q,p$,
$-\ln x\geq1-x$ implies
\[
	\sum_aq_a\ln(q_a/p_a)
	\geq\sum_a(q_a-p_a)=0,
\]
with equality only if $q=p$. Taking $p$ to be the product of
the marginals proves strict conditional mutual information.
Thus the literal four-ring fails SAL.

Its non-ZCL claim remains correct. The two-bit probability of
$00$ is $14/41$ on the four-ring, but on the three-ring it is
\[
	\frac{(W^2)_{00}}{\tr W^3}
	=\frac{5/4}{7/2}=\frac5{14}.
\]
They differ because $196\ne205$. Literal physical-trace
idempotence would identify the reduced closures and their
normalizations. A common rescaling of $W$ changes neither
probability nor entropy. This removes the printed finite-periodic
SAL witness; it does not disprove the independence of SAL and ZCL.

\section{A coefficient-separation argument for the fundamental theorem}\label{sec:matching}

The matching step in the proof of source Theorem~\ref{cpsv-thm1}
invokes independence of the $A$-family together with one $B$-sector,
although the proportionality relation contains all $B$-sectors.
The following coefficient-separation argument supplies this step and
preserves the conclusions of that theorem and Corollary~\ref{cpsv-II_cor2}.
Only active nonzero canonical sectors are retained: zero-weight blocks
carry no information about dimensions.

The elementary fact needed twice is the following. If
$z_1,\ldots,z_m$ are distinct and nonzero and
$\sum_ac_az_a^N=0$ for every $N\geq N_0$, use the $m$ equations
at $N_0,\ldots,N_0+m-1$. Their coefficient determinant is
\[
	\left(\prod_{a=1}^mz_a^{N_0}\right)
	\prod_{a<b}(z_b-z_a)\ne0.
\]
Thus every $c_a=0$. In particular, a nonempty power sum of
nonzero copy weights cannot vanish on an entire tail: after
equal weights are grouped, its coefficients are positive integer
multiplicities.

For the proportional theorem, form the union of the two BNT
families and identify representatives related by phase and gauge.
The remaining distinct classes have asymptotically orthogonal
normal MPVs by the source normal-tensor overlap lemma. Their
Gram matrix tends to the identity, so they are jointly linearly
independent at every sufficiently large length.

Suppose a sector of $A$ has no partner in $B$. In the identity
\[
	\ket{V^{(N)}(A)}=c_N\ket{V^{(N)}(B)},\qquad c_N\ne0,
\]
express all matched sectors using the chosen union representatives.
Coefficient comparison forces $\sum_q\mu_{j,q}^N=0$ for the
unmatched sector on the whole tail, contrary to the determinant
argument. Interchanging the two families excludes unmatched
sectors of $B$. This gives bijective phase-gauge matching without
a unit-weight assumption in each sector.

In the equal-MPV case, write the matched representatives as
$B_j^i=e^{i\phi_j}Y_jA_j^iY_j^{-1}$.
Coefficient comparison first gives only
\[
	\sum_q\mu_{j,q}^N
	=\sum_q(e^{i\phi_j}\nu_{j,q})^N
	\qquad(N\geq N_0).
\]
Group the distinct nonzero weights from both sides and subtract.
The same determinant makes every multiplicity difference zero.
The multisets therefore agree, their power sums agree at all
positive lengths, and matching the copies gives the global
similarity in the corollary.

For proportional families the common scalar $c_N$ remains in
every sector identity. The argument then establishes only the
sector matching asserted by the source theorem, not copy-weight recovery.

\section{Local corrections}\label{sec:local}

\subsection{CID requires positive complementary gaps}

The forward implication from transfer idempotence to the printed
CID definition has an independent adjacency obstruction.
Take the normal tensor
\[
	A^{(a,b)}=\frac1{\sqrt2}\ketbra{a}{b},
	\qquad a,b\in\{0,1\}.
\]
Its matrices span $M_2$, and
\begin{align*}
	\mc E_A(H)
	             & =\frac12\sum_{a,b}H_{bb}\ketbra{a}{a}
	=\frac12\tr(H)\Id_2,                                 \\
	\mc E_A^2(H) & =\frac12\tr(H)\Id_2.
\end{align*}
It is a single normal RFP, with no copy weights. Each physical
site contains two qubits, and its periodic MPS is the product
of Bell pairs between neighboring sites.

On a four-ring, place $Z$ on the right qubit of site one and
the left qubit of site two. These qubits share
$\ket{\varphi}=(\ket{00}+\ket{11})/\sqrt2$, so
\[
	\bra{\varphi}Z\otimes Z\ket{\varphi}
	=\tfrac12(1+1)=1.
\]
Shift the second observable to site three, as permitted by the
printed CID interval. The observables now act on different Bell
pairs, and the expectation is $[\frac12(1-1)]^2=0$.
Thus idempotence does not imply CID for arbitrary disjoint,
possibly adjacent regions.

The correct forward statement requires at least one free site
in each complementary gap, both before and after translation.
Every intervening power then satisfies $\mathbb E^g=\mathbb E$.
At adjacency, $g=0$ instead gives $\mathbb E^0=\Id$.
This is a substantive positive-gap correction, not a copy-weight
normalization.

\subsection{Square-root normalization}

In source Theorem~\ref{cpsv-thm:charact-MPS} and
Lemma~\ref{cpsv-lem:charact-NT-pure-RFP}, a trace-normalized
diagonal fixed point $\Lambda$ requires the factor
$\sqrt{\Lambda}$ rather than $\Lambda$.
The printed pair-index isometry condition gives
\begin{align*}
	\left[\sum_iU^iH(U^i)^\dagger\right]_{aa'}
	 & =\sum_{b,b'}H_{bb'}\delta_{aa'}\delta_{bb'} \\
	 & =\delta_{aa'}\tr H.
\end{align*}
Hence $A^i=\sqrt{\Lambda}U^i$ gives
$\mc E_A(H)=\Lambda\tr H$, whereas $A^i=\Lambda U^i$
gives $\Lambda^2\tr H$, of spectral radius
$\tr(\Lambda^2)$ rather than one.

Equivalently, one may retain $\Lambda$ as the Schmidt-amplitude
matrix but then require $\tr(\Lambda^2)=1$.
If $\Lambda$ denotes the density fixed point, the bond state in
the same theorem has amplitudes $\sqrt{\Lambda_{aa}}$.
This correction does not remove the independent copy-phase
obstruction.

\subsection{Indices, exponents, and endpoints}

The component identity in source Eq.~\eqref{cpsv-AA=A} should be
\[
	A^{i_1}A^{i_2}=\sum_jU_{(i_1,i_2),j}A^j.
\]
The indices $i_1,i_2$ remain free, and only $j$ is summed.
In ``Commuting parent Hamiltonians'', the length-$L$ support
vectors likewise require
\[
	\ket{v_{\alpha,\beta}}
	=\sum_{i_1,\ldots,i_L}
	(A^{i_1}\cdots A^{i_L})_{\alpha,\beta}
	\ket{i_1,\ldots,i_L}.
\]

In Appendix~C's proof of Theorem~\ref{cpsv-thm:IV.13}, the
identity
$A_\gamma=(m_\gamma/n_\gamma)
	U_\gamma M_\gamma U_\gamma^\dagger$
contributes $L$ scalar factors to a length-$L$ operator.
The first simultaneous blocked-operator representation must
therefore have coefficient
\[
	\left(\frac{m_\gamma}{n_\gamma}\right)^L\tr(\nu_\gamma^L),
\]
not $m_\gamma^L\tr(\nu_\gamma^L)/n_\gamma$.
This agrees with the later trace-power expression in that proof.

The entropy equality following the definition of SAL must read
\[
	S_L+S_{N-L}=S_{L+1}+S_{N-L-1}.
\]
The closing proof of Theorem~\ref{cpsv-thm:main-simple}
also cyclically interchanges the roles of conditions~(iii)
and~(iv) in its proposition assignments.
Proposition~\ref{cpsv-prop2to5} starts from condition~(ii).
Correcting these labels does not repair the false implications
established above.

Finally, the conjugating operator in source
Example~\ref{cpsv-ExZCLnoSAL} should act on the two different
qubits of a site:
\[
	U_n=\sigma_x^{(n,l)}\otimes\sigma_x^{(n,r)}.
\]
The exception list for strict failure of SAL must also include
$p=1$. At that endpoint the channel is deterministic on-site
unitary conjugation. It preserves the entropies of the Bell-pair
ring, which has two maximally mixed boundary qubits at each
proper connected cut: $S_L=2$, $S_N=0$, and $I_L=4$.
This endpoint correction leaves the non-SAL witness at $p=1/4$
unaffected.

% Keep the single reference in the current grid; the APS full-width
% bibliography separator otherwise forces an orphan page.
\renewcommand{\bibsection}{\par\medskip}
\bibliographystyle{apsrev4-2}
\bibliography{references}
\end{document}
